\documentclass[12pt,spanish]{article}
\usepackage[utf8]{inputenc}
\usepackage{babel}
\usepackage{listings}
\usepackage{mathpazo}
\usepackage{enumitem}
\usepackage{courier}
\usepackage{textcomp}
\usepackage{xcolor}
\usepackage{parskip}
\usepackage{fullpage}

\newcommand{\onelinerule}{\rule[2.3ex]{0pt}{0pt}}
\newcommand{\twolinerule}{\rule[6.2ex]{0pt}{0pt}}
\newcommand{\respuesta}{\framebox[\textwidth]{\twolinerule}}
\newcommand{\nombre}{%
  \begin{tikzpicture}[xscale=.4,yscale=.7]
    \draw (0, 0) rectangle (22, 1);
  \end{tikzpicture}%
}
%\newcommand{\rol}   {\framebox[0.3\textwidth]{\onelinerule}}
\newcommand{\rol}{%
  \begin{tikzpicture}[xscale=.4,yscale=.7]
    \draw[gray!40] ( 0, 0) grid      ( 9, 1);
    \draw          ( 0, 0) rectangle ( 9, 1);
    \draw          (10, 0) rectangle (11, 1);
    \draw (9 + .2, .5) -- (10 - .2, .5);
  \end{tikzpicture}%
}
\newcommand{\li}{\lstinline}
\providecommand{\pond}[1]{[{\small\textbf{#1\%}}]}

\lstdefinelanguage{py}{%
  classoffset=0,%
    morekeywords={%
      False,class,finally,is,return,None,continue,for,lambda,try,%
      True,def,from,nonlocal,while,and,del,global,not,with,print,%
      as,elif,if,or,yield,assert,else,import,pass,break,except,in,raise},%
    keywordstyle=\color{black!80}\bfseries,%
  classoffset=1,
    morekeywords={int,float,str,abs,len,raw_input,exit,range,min,max,%
      set,dict,tuple,list,bool,complex,round,sum,all,any,zip,map,filter,%
      sorted,reversed,dir,file,frozenset,open,%
      array,zeros,ones,arange,linspace,eye,diag,dot},
    keywordstyle=\color{black!50}\bfseries,%
  classoffset=0,%
  sensitive=true,%
  morecomment=[l]\#,%
  morestring=[b]',%
  morestring=[b]",%
  stringstyle=\em,%
}

\lstdefinelanguage{testcase}{%
  moredelim=[is][\bfseries]{`}{`},%
  backgroundcolor=\color{gray!20},%
}

\lstdefinelanguage{file}{%
  frame=single,%
}

\lstset{language=py}
\lstset{basicstyle=\ttfamily}
\lstset{columns=fixed}
\lstset{upquote=true}
\lstset{showstringspaces=false}
\lstset{rangeprefix=\#\ }
\lstset{includerangemarker=false}

\newlist{certamen}{enumerate}{1}
\setlist[certamen]{%
  label=\arabic*.,
  font=\LARGE\bfseries,%
  labelindent=-.5in,%
  leftmargin=0pt,%
  labelsep=1em%
}



\lstset{language=file,frame=single}

\begin{document}
  \thispagestyle{empty}
  \section*{Miércoles 7 de noviembre}

  El archivo \texttt{curso.txt} contiene los datos
  de los estudiantes de un curso.
  Cada línea corres\-ponde a un estudiante,
  y contiene varios datos separados por dos puntos («\verb+:+»)
  como en el siguiente ejemplo:
  \begin{verbatim}Martin:Soto:1990/8/24:39:86:22:56:30\end{verbatim}

  Los dos primeros datos son el nombre y el apellido.
  A continuación está la fecha de nacimiento
  en formato \texttt{año/mes/día}.
  Finalmente, están las cinco notas obtenidas por el estudiante.

  \begin{enumerate}[leftmargin=0pt]

    \item
      Escriba un programa llamado \texttt{clasificar.py} que,
      a partir de los datos de \texttt{curso.txt}, genere dos archivos
      llamados \texttt{aprobados.txt} y \texttt{reprobados.txt},
      que tengan los datos de los estudiantes aprobados y reprobados,
      respectivamente,
      usando el formato \texttt{nombre:apellido:promedio}.
      Por ejemplo,
      las primeras líneas de \texttt{aprobados.txt} serían éstas:
      \lstinputlisting[linerange=-5]{aprobados.txt}

    \item
      Escriba un programa que
      cree un nuevo archivo llamado \texttt{promedios.txt},
      que tenga los promedios y el estado de aprobación
      de los estudiantes.

      Las primeras líneas del archivo debe verse así:
      \lstinputlisting[linerange=-10]{promedios.txt}

      Note que los promedios están alineados a la derecha.
      Puede usar la siguiente plantilla para lograr esto
      usando el método \li!.format(...)!:
      \begin{lstlisting}[language=py,frame=none]
plantilla = '{0:20} {1:>4} {2}\n'
      \end{lstlisting}

      \newpage

    \item
      Escriba el programa \texttt{reporte.py}
      que muestre la siguiente información por pantalla:
      \lstinputlisting[language=testcase]{reporte.txt}

    \item
      Escriba el programa \texttt{pagina\_web.py}
      que genere el archivo \texttt{curso.html},
      cuyo contenido sea como el siguiente:
      \lstinputlisting[basicstyle=\small\ttfamily]{ejemplo.html}
      (Por brevedad aquí sólo se muestra los primeros dos estudiantes;
      usted debe agregarlos todos).
      Los datos listados son:
      el nombre completo,
      las notas de los cinco certámenes
      y el promedio final (encerrado entre \verb+<b>+ y \verb+</b>+).

      El archivo generado es una página web.
      Arrástrelo a su navegador favorito
      para verificar que fue creado correctamente.

  \end{enumerate}


\end{document}

